\scanpage{20}
\deleted{\unsure{Pour que}} \deleted{\unsure{Il contient l'unité} (i.e. est identique à $\mathcal A$)}
Pour que \inserted{$\underline b$ ait une unité, i.e. pour que} $\underline b=\mathcal A$, il faut et il suffit que $\underline a=\mathcal A$, i.e. que $T$ soit une \uline{trace finie} (i.e. $T(A)<\infty$ pour tout $A\in\mathcal A^+$).

\deleted{\uline{Proposition 1} Soit $A\in\mathcal A$ un élément hermitien de $\underline b$, \unsure{non nul}. Alors les projecteurs spectraux \inserted{$\unsure{E}$} sont dans $\underline a$.
\deleted{En effet, les projecteurs spectraux de $\underline b$ \unsure{sont dans}} \deleted{\unsure{voisin}}
Soit $E$ un projecteur spectral de $A$; il faut prouver $T(E)<+\infty$. \unsure{On}
}

\deleted{\uline{Proposition 1} Soit $A$ un opérateur \uline{normal} dans $\underline b$, \deleted{\unsure{un}} soit \deleted{$\unsure{\varphi}$} $f$ une fonction \deleted{\unsure{continue}} borélienne sur $\underline C$, \inserted{bornée, \deleted{à \unsure{support compact},}} \deleted{\unsure{support}} nulle au voisinage de $0$. Alors $f(A)\in\underline a$.

\deleted{\unsure{Admet} \illegible{4} \unsure{soit} \illegible{2}} \deleted{\unsure{Comme il} \illegible{3}}

\deleted{\uline{Corollaire 1}}

$A$ appartient à une sous-algèbre \unsure{uniformt} fermée auto-adjointe commutative $\underline c\subset\underline b$, $\underline c$ s'identifie à un $C_0(\Omega)$, $A$ à un élément $\varphi$ de $C_0(\Omega)$. $f(A)$ est l'opérateur défini par la fonction borélienne $f(\varphi)$ sur $C_0(\Omega)$. Cette dernière est \inserted{bornée} à support compact, donc \inserted{cette fonction est majorée,} majorée par $E$ fonction \deleted{caractéristique d'un} \inserted{continue à support} compact $K$. Il suffit donc de prouver le

\uline{Corollaire} \deleted{Soit} $\mathcal K(\Omega)\subset\underline a$
}
\mspage{21}
